\documentclass[letterpaper,11pt]{article}

\usepackage{latexsym}
\usepackage[empty]{fullpage}
\usepackage{titlesec}
\usepackage{marvosym}
\usepackage[usenames,dvipsnames]{color}
\usepackage{verbatim}
\usepackage{enumitem}
\usepackage[hidelinks]{hyperref}
\usepackage{fancyhdr}
\usepackage[english]{babel}
\usepackage{tabularx}
\input{glyphtounicode}

\pagestyle{fancy}
\fancyhf{}
\fancyfoot{}
\renewcommand{\headrulewidth}{0pt}
\renewcommand{\footrulewidth}{0pt}

\addtolength{\oddsidemargin}{-0.5in}
\addtolength{\evensidemargin}{-0.5in}
\addtolength{\textwidth}{1in}
\addtolength{\topmargin}{-.5in}
\addtolength{\textheight}{1.0in}

\urlstyle{same}

\raggedbottom
\raggedright
\setlength{\tabcolsep}{0in}

\titleformat{\section}{
  \vspace{-4pt}\scshape\raggedright\large
}{}{0em}{}[\color{black}\titlerule \vspace{-5pt}]

\pdfgentounicode=1

\newcommand{\resumeItem}[1]{
  \item\small{
    {#1 \vspace{-2pt}}
  }
}

\newcommand{\resumeSubheading}[4]{
  \vspace{-2pt}\item
    \begin{tabular*}{0.97\textwidth}[t]{l@{\extracolsep{\fill}}r}
      \textbf{#1} & \textit{\small #4} \\
      \textit{\small#3} & \textit{\small #2} \\
    \end{tabular*}\vspace{-7pt}
}

\newcommand{\resumeProjectHeading}[2]{
    \item
    \begin{tabular*}{0.97\textwidth}{l@{\extracolsep{\fill}}r}
      \small#1 & #2 \\
    \end{tabular*}\vspace{-7pt}
}

\renewcommand{\labelitemi}{$\bullet$}
\renewcommand\labelitemii{$\bullet$}

\newcommand{\resumeSubHeadingListStart}{\begin{itemize}[leftmargin=0.15in, label={}]}
\newcommand{\resumeSubHeadingListEnd}{\end{itemize}}
\newcommand{\resumeItemListStart}{\begin{itemize}[leftmargin=0.2in]}
\newcommand{\resumeItemListEnd}{\end{itemize}\vspace{-5pt}}

\begin{document}

\begin{center}
    \textbf{\Huge \scshape Jodh Singh Nahal} \\ \vspace{1pt}
    \small 778-538-2351 $|$
    \href{mailto:nahal.jodhsingh@gmail.com}{\underline{nahal.jodhsingh@gmail.com}} $|$
    \href{https://www.linkedin.com/in/jodhsinghnahal/}{\underline{linkedin.com/in/jodhsinghnahal}} $|$
    \href{https://github.com/jodhsinghnahal}{\underline{github.com/jodhsinghnahal}}
\end{center}

\section{Education}
\resumeSubHeadingListStart
    \resumeSubheading
      {University of British Columbia}{Expected May 2028}
      {Bachelor of Applied Science in Computer Engineering}{Vancouver, BC}
      \resumeItemListStart
        \resumeItem{Presidential Scholars Award (\$60,000), Dean's Honour List, available 4 or 8 months, CGPA: 84.7\%.}
      \resumeItemListEnd
\resumeSubHeadingListEnd

\section{Technical Skills}
\begin{itemize}[leftmargin=0.15in, label={}]
\small{\item{
     \textbf{Protocols}{: CAN/CAN-FD, MQTT, Bluetooth, UART, SPI, I2C, Modbus, NMEA 2000, RV-C, PCIe} \\
     \textbf{Languages}{: C/C++, Python, Shell, SQL} \\
     \textbf{Tools}{: Linux/Unix, Logic Analyzers, Oscilloscopes, BLE Sniffers, PCAN-USB, Git}
}}
\end{itemize}

\section{Experience}
\resumeSubHeadingListStart

    \resumeSubheading
      {Software Test / Embedded Development Co-op}{May 2024 -- Aug 2025}
      {Xantrex}{Burnaby, BC}
      \resumeItemListStart
        \resumeItem{Developed an object-oriented, multi-threaded, non-blocking C++ firmware architecture on Linux-based Raspberry Pi gateways to parse RV-C CAN frames from STM32 devices and publish structured messages over MQTT.}
        \resumeItem{Extended firmware support to NMEA 2000 and a proprietary UART battery communication protocol.}
        \resumeItem{Troubleshot CAN/BLE communication faults on embedded inverter and charger platforms using oscilloscopes, logic analyzers, PCAN-USB adapters, and BLE sniffers.}
        \resumeItem{Built Python Hardware-in-the-Loop (HIL) CAN and UART simulators to run 100+ verification tests on the Raspberry Pi gateway, validating edge cases and stress-testing real-time performance.}
        \resumeItem{Achieved $7\times$ device capacity, 90\% lower CPU usage, and $\sim$2\,s lower latency after the C++ migration.}
        \resumeItem{Replaced a $\sim$60k-line Python implementation with the C++ architecture.}
      \resumeItemListEnd

    \resumeSubheading
      {Hardware Verification Engineering Intern}{Sept 2026 -- Present}
      {AMD}{Vancouver, BC}
      \resumeItemListStart
        \resumeItem{Contributing to verification of the IOHUB subsystem for next-generation AMD processor architectures.}
        \resumeItem{Verifying IOMMU L1 cache behavior using UVM in a constrained-random testbench.}
        \resumeItem{Expanded the constrained-random testbench to cover corner cases and improve functional coverage.}
      \resumeItemListEnd

    \resumeSubheading
      {Firmware / Embedded Systems Developer}{Sept 2023 -- July 2026}
      {UBC Rover}{Vancouver, BC}
      \resumeItemListStart
        \resumeItem{Led migration to a CAN-FD architecture using Moteus r4.11 controllers and BLDC motors for closed-loop torque control and higher bandwidth.}
        \resumeItem{Interfaced an Intel NUC (onboard PC) with SocketCAN and bridged data to ROS.}
        \resumeItem{Led a team of 3 engineers building a distributed embedded control system across ESP32 (FreeRTOS), STM32, Arduino, and TI LaunchPad MCUs over SPI/I2C.}
        \resumeItem{Aggregated inputs from switches, buttons, potentiometers, encoders, and joysticks and streamed data over UART.}
        \resumeItem{Developed a 16-channel telemetry system on a TI C2000 (LAUNCHXL-F28069M) using ePWM-triggered ADC interrupts, streaming GPIO/ADC data over SCI/UART.}
      \resumeItemListEnd

    \resumeSubheading
      {Undergraduate Research Assistant}{May 2026 -- Sept 2026}
      {UBC Electrical and Computer Engineering --- Supervisor: Dr. Lutz Lampe}{Vancouver, BC}
      \resumeItemListStart
        \resumeItem{Researching battery-less platforms for the Internet of Underwater Things via simulation of piezoelectric interface circuits.}
        \resumeItem{Optimizing power extraction of piezoelectric energy harvesting interface circuits (SEH, SSHI, SECE) through LTspice and Simulink/Simscape simulation.}
      \resumeItemListEnd

\resumeSubHeadingListEnd

\section{Projects}
\resumeSubHeadingListStart

    \resumeProjectHeading
      {\textbf{OnlyVim} $|$ \emph{Python, Linux, D-Bus, RPC}}{}
      \resumeItemListStart
        \resumeItem{Developing a Linux daemon enabling global text editing with Neovim keybindings via keyboard remapping, D-Bus, and RPC interfaces.}
        \resumeItem{Designed low-latency event handling and Linux IPC for system-wide integration across applications.}
      \resumeItemListEnd

    \resumeProjectHeading
      {\textbf{Multimodal AI Messenger} $|$ \emph{Node.js, Express, WebSocket, MongoDB, Gemini API}}{}
      \resumeItemListStart
        \resumeItem{Built a concurrent full-stack messaging platform with Node.js, MongoDB, and WebSockets for real-time updates.}
      \resumeItemListEnd

\resumeSubHeadingListEnd

\end{document}
